\documentclass[11pt]{article}
\usepackage{amsmath,amssymb,amsthm}
\newtheorem{theorem}{Theorem}
\newtheorem{lemma}{Lemma}
\title{Cauchy Reflection and Monotone Readout Coordinates}
\author{}\date{}
\begin{document}\maketitle
\paragraph{Status.}
Complete written algebraic identities with exact rational checks, not Lean.
The monotone-readout feasibility property below remains open as the number
of butterfly levels grows. These identities are not a compiler theorem.

Let $p_1,\ldots,p_H$ be distinct real numbers with $p_i<T/2$ and define
\[
 C_{ij}=\frac1{T-p_i-p_j},\qquad
 Q(x)=\prod_i(x-p_i),\qquad A(x)=Q(T-x),\qquad
 w_i=\frac{A(p_i)}{Q'(p_i)},\qquad B=C\operatorname{diag}(w).
\]
\begin{theorem}[Exact reflection]
For a vector $v$, let $N_v$ be the unique polynomial of degree less than
$H$ with $N_v(p_i)=A(p_i)v_i$. Then
\[
 (Bv)_i=\frac{N_v(T-p_i)}{A(p_i)},\qquad B^2=I,
 \qquad C^{-1}=\operatorname{diag}(w)C\operatorname{diag}(w).
\]
For every diagonal multiplier $D_u$,
\[
 C D_u C^{-1}=B D_u B.
\]
\end{theorem}
\begin{proof}
The partial fraction identity for $N_v/Q$ is
\[
 \frac{N_v(x)}{Q(x)}
 =\sum_j\frac{A(p_j)v_j}{Q'(p_j)(x-p_j)}.
\]
Evaluate at $x=T-p_i$. The interpolation polynomial associated with
$Bv$ is therefore exactly $N_v(T-x)$, whose degree is still less than
$H$. Applying reflection twice proves $B^2=I$. Since every $w_i$ is
nonzero, this proves the inverse formula. Diagonal matrices commute,
so the last identity follows as well.
\end{proof}

\begin{lemma}[Exact symmetric Schur complement]
Remove an index set $J$ from the positive definite Cauchy matrix. Its
Schur complement on the remaining indices is
\[
 C^J_{ik}=f_J(p_i)C_{ik}f_J(p_k),\qquad
 f_J(x)=\prod_{j\in J}\frac{x-p_j}{T-x-p_j}.
\]
\end{lemma}
\begin{proof}
For one removed index $j$, direct subtraction gives
\[
 C_{ik}-\frac{C_{ij}C_{jk}}{C_{jj}}
 =\frac{p_i-p_j}{T-p_i-p_j}\,C_{ik}\,
   \frac{p_k-p_j}{T-p_k-p_j}.
\]
Iterating gives the displayed product. The diagonal factors are nonzero
on remaining distinct poles. Positive definiteness follows either from
Schur complements or from
\[
 v^\top C v=\int_0^\infty e^{-Tt}
       \left(\sum_i v_i e^{p_it}\right)^2dt>0
\]
for nonzero $v$, since distinct real exponentials are linearly independent.
The integral converges under $2\max_i p_i<T$.
\end{proof}

\paragraph{Application to the signed Haar decoder.}
For a fixed routing, let $u_g(y)$ be its signed Haar block mask in output
coordinates. If $x=Cc$, then every unbiassed block readout at a grid query
$p_s$ is
\[
 (C D_{u_g}c)_s=(B D_{u_g}B x)_s.
\]
This changes no scores and imposes no restriction on signed residues.
The family of all multiplier operators commutes, and each indicator
multiplier becomes a projection in these coordinates. This does not
make its matrix entries nonnegative.

Order the poles increasingly. A strictly increasing readout vector can
be written uniquely as
\[
 x=\alpha\mathbf1+\sum_{j=1}^{H-1}d_j s_j,
 \qquad d_j>0,
\]
where $s_j$ is the indicator of the final $H-j$ ordered coordinates.
Eliminating the independent block biases by comparing each required
positive query with each required negative query gives a real matrix
$L$ in residue coordinates. The exact strengthened feasibility question
is
\[
 L C^{-1}\left(\alpha\mathbf1+\sum_jd_js_j\right)>0,
 \qquad d_j>0.
\]
It is stronger than the original signed biased decoder. No theorem here
asserts it for all routings or target words. A fixed common increasing
$x$ is stronger still and is not assumed.

Exact finite probes show that the physical-alternation stress word at
$H=32$ has a feasible strictly increasing readout, and the same holds for
the tested $H=16$ words. A greedy construction by successively choosing
one dominating gap fails already for the $H=8$ stress word: every gap
column has a negative entry among the required comparison rows, and
neither sign of the constant column is nonnegative on all these rows.
Thus these examples require balanced cancellations in the gap variables.

Finite feasibility rates alone provide little asymptotic evidence.
The signed biased family has $2H-1$ parameters and
$q=H\log_2(H)/2$ queries. For a generic rank-$d$ central linear image in
$\mathbb R^q$, the fraction of orthants met is
$2^{1-q}\sum_{j=0}^{d-1}\binom{q-1}{j}$.
For $d=2H-1$, this is already close to one at the small sizes tested.
The desired growing-capacity statement must use the coupling between
routing and target signs, rather than this dimension count.

\paragraph{Limits of a positivity shortcut.}
The paired divided-difference matrix
$K_{ij}=1/((t_j-r_i)(t_j-s_i))$ is not universally sign regular for
arbitrary matching of two separated pole clusters. For four points,
match a sorted set $R$ to its translate by $100$ using the permutation
$(0,3,1,2)$. The exact determinant is negative for
\[
 R=(0,14,17,31),\qquad t=(229,245,338,379),
\]
and positive for
\[
 R=(0,8,13,21),\qquad t=(207,229,257,274).
\]
Both pole sets are symmetric binary grids and all query points are
exterior. Thus a proposed recursive positivity argument must use
additional hierarchy or routing information. The symmetric Schur
identity above remains valid; it does not identify these arbitrary
paired matrices with symmetric Schur complements.
\end{document}
